\documentclass{article}
\usepackage[margin=2cm]{geometry}
\usepackage{amsmath}
\usepackage{graphicx}
\usepackage{booktabs}
\usepackage[utf8]{inputenc}
\usepackage{ragged2e}
\setlength{\emergencystretch}{3em}
\begin{document}
\sloppy

which for the coset representative is given by


In addition to performing the above Wick rotation, we also perform a Wick rotation on the.
sigma model


\begin{equation}
{ \begin{array} { c c c c c c } { { \frac { S O ( 4, 1 ) } { S O ( 4 ) } } \rightarrow { \frac { S O ( 4, 1 ) } { S O ( 3, 1 ) } } \simeq d S _ { 4 } ~ ~ } & { } & { } & { } & { } & { } & { } & { ( \mathbf { 2 } ) } \\ { \dots } & { } & {. } & { } & { } & { } & { } & { } & { ( \mathbf { 2 } ) } \end{array} }
\end{equation}


The metric on the sigma model is now that of dS4, as opposed to the H4 that we had in the
Lorentzian case . This can be obtained by making the following change to the H4 coset,


\begin{equation}
\phi _ { r } \rightarrow i \phi _ { r } ~ r = 1, 2, 3 ~
\end{equation}


It would be interesting to understand this analytic continuation from first principles and its.
relation to Euclidean supersymmetry, possibly along the lines of . For now, we just note that
such a Wick rotated model seems necessary to obtain regular, supersymmetric solutions.


\subsection*{0.1Euclidean supersymmetry}


The next task is to identify the form of the Euclidean supersymmetry variations. Motivation
for the form of these variations may be obtained by analysis of the free differential algebra
(FDA) of the F(4) gauged supergravity theory with Hg vacuum, as discussed in Appendix .
The final result for this FDA is given in, and is noted to be of the same form as the FDA
for the theory with dS6 background (identified in ), with two differences. The first obvious
difference is that the metrics differthe space considered in was dS6 with mostly minus
signature, whereas we are currently focused on positive definite H6. However, both of these
spaces have Rv = 2Om2gv. The second difference is in the definition of Dirac conjugate
spinors. However, once the difference in definition of the gamma matrices is accounted for,.
the only difference is a factor of i, i.e.


Because of these similarities, the supersymmetry variations in the current case are expected
to be of a similar form to that of . In particular, the variations of the fermions are expected
to be of the form


\begin{equation}
\begin{array} { c } { \delta \chi _ { A } = - \frac { 1 } { 2 } \gamma ^ { \mu } \partial _ { \mu } \sigma \varepsilon _ { A } + N _ { A B } \varepsilon ^ { B } + \dots } \\ { \delta \psi _ { A \mu } = D _ { \mu } \varepsilon _ { A } + i S _ { A B } \gamma _ { a } \varepsilon ^ { B } + \dots } \\ { \delta \lambda _ { A } ^ { I } = - \hat { P } _ { n l } ^ { I } \sigma _ { A l } ^ { \tau } \partial _ { \mu } \phi ^ { j } \varepsilon ^ { B } + \frac { 1 }}\end{array}
\end{equation}

\end{document}
